% Generated from validated Article Draft Result. Do not hand-edit; revise the structured draft instead.
\section{汎用Agent競争の外側 — Deep Thinkと科学推論}
\label{pkg:specialized-reasoning-science}
\noindent\textbf{coding agentが2月の主旋律だった一方、Google DeepMindはGemini 3 Deep Thinkをscience・research・engineering向けのspecialized reasoning modeとして更新した。汎用modelの一律な能力向上ではなく、難しい推論workloadへ計算と利用形態を寄せる別のfrontierが見えている。}\autocite{src-dd5d9b8662f3}

Google DeepMindは2月12日、Gemini 3 Deep Thinkの更新をscience、research、engineering problem向けのspecialized reasoning modeとして発表した。この位置付けは重要だ。Deep Thinkの結果をGemini 3 family全体の標準能力として一般化するのではなく、特定の難度・領域へ推論資源を寄せたmodeとして読む必要がある。2月のfrontier競争をagentだけで説明すると、この『用途を絞って推論を深くする』方向を見落とす。\autocite{src-dd5d9b8662f3}


\begin{claimboundary}
発表時点のavailabilityにも境界がある。Deep ThinkはGemini appでGoogle AI Ultra subscriber向けに提供され、APIについてはearly-access pathとして説明された。したがって、2月時点で一般的なGemini API modelと同じ広いavailabilityを持っていたとは書けない。specialized reasoningの技術的意味と、誰がどのsurfaceから利用できたかは分けて扱う。\autocite{src-dd5d9b8662f3}
\end{claimboundary}

\begin{claimboundary}
科学・工学領域のbenchmarkや他modelとの比較結果はGoogle側のvendor claimであり、本稿で独立再現したものではない。Deep Thinkがscience/research/engineering向けに更新されたことは一次資料で固定できるが、そこで示されたscoreをもって一般的な『科学能力の優位』を確定することはしない。specialized modeという設計上の事実と、performance superiorityの主張を分離する。\autocite{src-dd5d9b8662f3}
\end{claimboundary}

同じ月のGemini APIには、2月19日にGemini 3.1 Pro Previewとcustom-tools endpointが記録されている。こちらはreasoningとtool useを一般的なAPI surfaceへ接続する更新であり、Deep Thinkとは役割が異なる。Deep ThinkをGemini 3.1 Pro Previewの上位版として単純に並べるより、前者はspecialized science/research reasoning、後者はpreview-stageのgeneral reasoning/tool-use APIとして別々に読む方が、2月のGoogle stackの広がりを正確に表せる。\autocite{src-a24fd97eb3d0}


\begin{claimboundary}
Gemini 3.1 Proは2月時点でPreviewであり、Gemini API changelogはidentifierとchronologyを確認する一次資料ではあっても、qualityやbenchmarkを独立検証する資料ではない。同じEvidence taskには2月26日のFlash Image Previewも含まれるが、本章ではimage generationを主題化しない。reasoning/tool-use側のPro PreviewだけをDeep Thinkとの役割比較に使い、別modalの更新はinteraction章へ委ねる。\autocite{src-a24fd97eb3d0}
\end{claimboundary}

2月の構図にDeep Thinkを残す意義は、frontierの評価軸が一つではないことを示す点にある。agentic engineeringは長時間のtool-using executionを伸ばす。一方、Deep Thinkはscience・research・engineeringという難しいproblem classへreasoning modeを特化する。3月以降を見る際にも、汎用modelのscoreだけでなく、『どのworkloadに対してspecialized modeやaccess pathが設計されるか』を独立した軸として追う必要がある。これは2月以降の具体的な成果を先取りする主張ではなく、2月時点で成立した技術的な問いである。\autocite{src-a24fd97eb3d0,src-dd5d9b8662f3}
